\documentclass[a4paper,12pt]{article}
\usepackage[T2A]{fontenc}
\usepackage[utf8]{inputenc}
\usepackage[russian]{babel}
\usepackage{amsmath,amssymb,mathtools}
\renewcommand{\familydefault}{\sfdefault}
\usepackage{icomma}
\usepackage[left=23mm,right=22mm,top=23mm,bottom=24mm]{geometry}
\usepackage{microtype}
\usepackage{xcolor}
\usepackage{tikz}
\usetikzlibrary{calc,arrows.meta,positioning,shapes.geometric}
\usepackage{tabularx,booktabs}
\usepackage{enumitem}
\usepackage{parskip}
\usepackage{fancyhdr}
\usepackage{setspace}
\usepackage{needspace}
\usepackage[hidelinks]{hyperref}
\definecolor{accent}{HTML}{087B80}
\definecolor{darkaccent}{HTML}{15575A}
\definecolor{textgray}{HTML}{303D43}
\definecolor{muted}{HTML}{62777D}
\color{textgray}
\setstretch{1.16}
\setlist[enumerate]{leftmargin=*,itemsep=.62em,topsep=.4em}
\setlist[itemize]{leftmargin=*,itemsep=.2em,topsep=.25em}
\newlist{checklist}{itemize}{1}
\setlist[checklist]{label=$\square$,leftmargin=*,itemsep=.35em,topsep=.3em}
\newcommand{\key}[1]{\textcolor{darkaccent}{\bfseries #1}}
\newcommand{\ruleline}{\par\vspace{.15em}{\color{accent!65}\hrule height .6pt}\vspace{.65em}}
\newcommand{\checks}[1]{\item[$\square$] #1}
\pagestyle{fancy}
\fancyhf{}
\fancyhead[L]{\footnotesize\color{muted}Чек-лист · Григорий Архипов}
\fancyhead[R]{\footnotesize\color{muted}21.10.26}
\fancyfoot[C]{\footnotesize\color{muted}\thepage}
\renewcommand{\headrulewidth}{0pt}
\setlength{\headheight}{15pt}

\begin{document}
\begin{titlepage}
\thispagestyle{empty}
\begin{center}
\vspace*{38mm}
{\color{accent}\small МАТЕМАТИКА · ПОДГОТОВКА К ЕГЭ}\par
\vspace{18mm}
{\fontsize{34}{42}\selectfont\bfseries Чек-лист}\par
\vspace{11mm}
{\fontsize{21}{29}\selectfont\bfseries Текстовые задачи\\на движение}\par
\vspace{5mm}
{\large Река · Средняя скорость · Протяжённые тела}\par
\vspace{22mm}
{\color{muted}\large 21 октября 2026 года}\par
\vspace{7mm}
{\color{darkaccent}\large Григорий Архипов}\par
\vfill
{\small Лёвин Артём Александрович}\par
{\small эксклюзивно для Григория Архипова}\par
\vspace{27mm}
\end{center}
\begin{tikzpicture}[remember picture,overlay]
\draw[accent!55,line width=1.5pt]
 ([xshift=22mm,yshift=20mm]current page.south west)
 .. controls ([xshift=56mm,yshift=32mm]current page.south west)
 and ([xshift=-59mm,yshift=3mm]current page.south east)
 .. ([xshift=-22mm,yshift=20mm]current page.south east);
\end{tikzpicture}
\end{titlepage}

\section*{01. Общий метод: от текста к уравнению}
\ruleline
\begin{checklist}
\checks{Выделяю \textbf{то, что требуется найти}; по возможности обозначаю именно эту величину через $x$.}
\checks{Заполняю таблицу «скорость --- путь --- время» и вычисляю формульный столбец по $S=vt$.}
\checks{Различаю \textbf{время движения}, \textbf{время стоянки} и \textbf{время от отправления до возвращения}.}
\checks{Составляю уравнение из смыслового условия: сумма, разность или отношение времён и путей.}
\checks{Проверяю допустимые значения, единицы и реалистичность ответа.}
\end{checklist}

\begin{center}
\renewcommand{\arraystretch}{1.35}
\begin{tabularx}{\linewidth}{>{\bfseries}lXXX}
\toprule
Этап & Скорость $v$ & Путь $S$ & Время $t$\\
\midrule
Первый & $v_1$ & $S_1$ & $S_1/v_1$\\
Второй & $v_2$ & $S_2$ & $S_2/v_2$\\
\bottomrule
\end{tabularx}
\end{center}

\key{Главный принцип.} Не пытайтесь находить каждую неизвестную величину. Если в вопросе требуется отношение $x/y$ или разность $v_1-v_2$, достаточно определить именно её.

\section*{02. Движение по реке}
\ruleline
Пусть $x$ км/ч --- собственная скорость судна, $u$ км/ч --- скорость течения. Тогда
\[
 v_{\text{по}}=x+u,\qquad v_{\text{против}}=x-u,\qquad
 v_{\text{плота}}=u.
\]
\key{Ограничения:} при движении судна против течения $x>u>0$; все времена должны быть положительными.

\begin{checklist}
\checks{По течению прибавляю скорость течения, против течения --- вычитаю.}
\checks{В строке стоянки записываю $S=0$, $v=0$, а время --- по условию; формулу $t=S/v$ для стоянки не применяю.}
\checks{Если участники отправились в разные моменты, учитываю сдвиг начала движения.}
\checks{Если надо найти отношение скоростей, сначала выражаю его через уравнение, не решая систему без необходимости.}
\end{checklist}

\subsection*{Пример 1. Разность времени}
Судно проходит 24 км против течения и столько же по течению. Скорость течения 2 км/ч. Против течения оно движется на 3 часа дольше. Найдём собственную скорость $x$.
\[
 \frac{24}{x-2}-\frac{24}{x+2}=3,\quad x>2.
\]
\[
 \frac{96}{x^2-4}=3\ \Longrightarrow\ x^2=36
 \ \Longrightarrow\ \boxed{x=6\text{ км/ч}}.
\]
Проверка: $24/4=6$ ч против течения, $24/8=3$ ч по течению; разность равна 3 ч.

\subsection*{Пример 2. Стоянка включается в общее время}
Теплоход вышел в 9:00, прошёл 48 км по течению, простоял 2 часа и вернулся по тому же маршруту в 20:00. Скорость течения --- 4 км/ч. Найдём собственную скорость $x$.
\[
 \underbrace{\frac{48}{x+4}}_{\text{по течению}}+
 \underbrace{2}_{\text{стоянка}}+
 \underbrace{\frac{48}{x-4}}_{\text{против течения}}
 =11,\quad x>4.
\]
\[
 \frac{96x}{x^2-16}=9\ \Longrightarrow\
 3x^2-32x-48=0\ \Longrightarrow\ \boxed{x=12\text{ км/ч}}.
\]
Проверка: $3+2+6=11$ ч. Отрицательный корень $-\tfrac43$ не удовлетворяет условию $x>4$.

\clearpage
\section*{03. Плот, отношения величин, пропорции}
\ruleline
\begin{checklist}
\checks{Скорость плота совпадает со скоростью течения.}
\checks{Если события завершились одновременно, сравниваю \textbf{все прошедшие интервалы времени}, учитывая разный старт.}
\checks{Если пути неизвестны, обозначаю общий ненулевой путь через $S$ и сокращаю его только после обоснования $S\ne0$.}
\checks{При преобразовании пропорции проверяю, что все используемые знаменатели ненулевые.}
\end{checklist}

\key{Пример 3. Одновременное прибытие при разном старте.} Плот со скоростью 4 км/ч прошёл 36 км за $36/4=9$ ч. Лодка отправилась часом позже, прошла 60 км по течению и 60 км обратно и вернулась одновременно с прибытием плота. Тогда
\[
 \frac{60}{x+4}+\frac{60}{x-4}=9-1=8,\quad x>4,
\]
\[
 \frac{120x}{x^2-16}=8\ \Longrightarrow\
 x^2-15x-16=0\ \Longrightarrow\ \boxed{x=16\text{ км/ч}}.
\]
Проверка: $60/20+60/12=3+5=8$ ч.

\key{Пример 4. Ищем отношение.} На одинаковые ненулевые расстояния $S$ баржа тратит по течению вдвое меньше времени, чем против течения. Пусть $x$ --- её собственная скорость, $u$ --- скорость течения:
\[
 \frac{S}{x+u}=\frac12\,\frac{S}{x-u},\quad
 S>0,\quad x>u>0.
\]
Сократим $S$ и перемножим крест-накрест:
\[
 2(x-u)=x+u\ \Longrightarrow\ x=3u
 \ \Longrightarrow\ \boxed{\frac{x}{u}=3}.
\]
\key{Смысл:} собственная скорость баржи в 3 раза больше скорости течения. Отдельно искать $x$ и $u$ не требуется.

\subsection*{Безопасные преобразования пропорций}
Если $a,b,c,d\ne0$, то из $\dfrac ab=\dfrac cd$ следует
\[
 ad=bc,\qquad \frac ac=\frac bd,\qquad
 \frac ba=\frac dc.
\]
Обратное преобразование дроби требует ненулевого числителя. Безопасная основа всех перестановок --- равенство произведений по диагонали.

\section*{04. Средняя скорость}
\ruleline
\begin{checklist}
\checks{Понимаю смысл: средняя скорость --- постоянная скорость для того же \textbf{общего пути} за то же \textbf{общее время}.}
\checks{Всегда начинаю с формулы $v_{\text{ср}}=\dfrac{S_1+S_2}{t_1+t_2}$.}
\checks{Отличаю половину \textbf{времени} от половины \textbf{пути}: это разные модели.}
\checks{При сложных дробях делю на дробь умножением на обратную и сокращаю множители до вычисления произведений.}
\end{checklist}

\key{Пример 5. Равные промежутки времени.} Скорости 84 и 56 км/ч, время на каждом участке равно $t>0$:
\[
 v_{\text{ср}}=\frac{84t+56t}{t+t}=\boxed{70\text{ км/ч}}.
\]
\key{Пример 6. Равные участки пути.} Те же скорости, но длина каждого участка равна $S>0$:
\[
 v_{\text{ср}}=\frac{2S}{S/84+S/56}
 =\frac{2\cdot84\cdot56}{84+56}=\boxed{67{,}2\text{ км/ч}}.
\]
\key{Ошибка:} число $70$ здесь получилось бы при усреднении скоростей без учёта времени. При равных путях более медленный участок занимает больше времени.

\clearpage
\section*{05. Движение протяжённых тел}
\ruleline
\begin{checklist}
\checks{Различаю точечный ориентир (столб, наблюдатель) и протяжённый объект (поезд, мост, лесополоса).}
\checks{При полном прохождении поезда мимо моста складываю длины поезда и моста.}
\checks{При прохождении поезда мимо одного наблюдателя учитываю \textbf{только длину проходящего поезда}.}
\checks{Для встречного движения беру сумму скоростей, для попутного обгона --- положительную разность.}
\checks{Начальный промежуток между объектами включаю в относительный путь, если он указан.}
\checks{Согласую единицы: м перевожу в км, секунды и минуты --- в часы, когда скорость дана в км/ч.}
\end{checklist}

\key{Идея относительной скорости.} Если два тела движутся вдоль одной прямой, их относительная скорость по модулю равна
\[
 v_{\text{отн}}=\begin{cases}
 v_1+v_2,&\text{навстречу друг другу},\\
 |v_1-v_2|,&\text{в одном направлении}.
 \end{cases}
\]
Время полного прохождения $t=S_{\text{отн}}/v_{\text{отн}}$; относительный путь определяется конкретным начальным и конечным положением тел.

\begin{center}
\renewcommand{\arraystretch}{1.32}
\begin{tabularx}{\linewidth}{>{\raggedright\arraybackslash}X>{\raggedright\arraybackslash}X}
\toprule
Ситуация & Относительный путь при полном прохождении\\
\midrule
Поезд мимо столба или неподвижного человека & длина поезда $L$\\
Поезд целиком проходит мост длиной $B$ & $L+B$\\
Один поезд полностью проходит мимо другого & $L_1+L_2$\\
Полный обгон с начальным зазором $d$ & $L_1+L_2+d$\\
\bottomrule
\end{tabularx}
\end{center}

\key{Пример 7. Мимо столба.} Поезд со скоростью 60 км/ч проходит мимо столба за 30 с. Его длина:
\[
 t=\frac{30}{3600}=\frac1{120}\text{ ч},\quad
 L=60\cdot\frac1{120}=0{,}5\text{ км}=\boxed{500\text{ м}}.
\]

\key{Пример 8. Мимо лесополосы.} Поезд со скоростью 90 км/ч полностью проходит лесополосу длиной 800 м за 1 мин. Длина поезда:
\[
 L+0{,}8=90\cdot\frac1{60}=1{,}5
 \quad\Longrightarrow\quad \boxed{L=700\text{ м}}.
\]

\key{Особое внимание.} Если пассажир \emph{внутри} движущегося поезда засекает, за сколько мимо него прошёл товарный поезд, пассажир служит точечным ориентиром. Длина пассажирского поезда в относительный путь не входит.

\section*{06. Расчёты без лишней арифметики}
\ruleline
\begin{checklist}
\checks{Перевожу $30$ с в $\tfrac1{120}$ ч, $1$ мин в $\tfrac1{60}$ ч, $600$ м в $0{,}6$ км.}
\checks{Дробь в знаменателе обращаю: $\dfrac{a}{b/c}=\dfrac{ac}{b}$ при $b\ne0$, $c\ne0$.}
\checks{В произведениях сначала сокращаю общие множители; при необходимости раскладываю числа на простые.}
\checks{После расчёта отвечаю ровно на тот вопрос, который сформулирован в задаче.}
\end{checklist}

\clearpage
\thispagestyle{fancy}
\begin{center}
{\LARGE\bfseries\color{darkaccent}Алгоритм: задача на движение}\par
\vspace{5mm}
{\small\color{muted}Выберите модель, составьте уравнение, проверьте ответ}\par
\vspace{9mm}
\end{center}
\tikzset{
 flow/.style={draw=darkaccent!75,rounded corners=1.4mm,align=center,text width=11cm,minimum height=11mm,inner sep=3mm,line width=.6pt,font=\small},
 choice/.style={diamond,draw=darkaccent!75,aspect=2.4,align=center,text width=2.5cm,inner sep=1mm,font=\small,line width=.6pt},
 branch/.style={draw=darkaccent!75,rounded corners=1.4mm,align=center,text width=3.5cm,minimum height=27mm,inner sep=2mm,line width=.6pt,font=\small},
 arr/.style={-Latex,draw=accent,line width=.85pt}
}
\begin{center}
\begin{tikzpicture}
\node[flow] (a) at (0,0) {\textbf{Прочитайте задачу}\quad Что известно и какую величину требуется найти?};
\node[choice] (b) at (0,-2.35) {Какая модель?};
\node[branch] (c1) at (-5,-5.1) {\textbf{Река}\\$v=x\pm u$\\Таблица пути и времени, учёт стоянки};
\node[branch] (c2) at (0,-5.1) {\textbf{Средняя скорость}\\$v_{\text{ср}}=S_{\text{общ}}/t_{\text{общ}}$\\Равные времена или пути?};
\node[branch] (c3) at (5,-5.1) {\textbf{Протяжённые тела}\\Относительная скорость, длины и зазоры};
\coordinate (m) at (0,-7.05);
\node[flow] (d) at (0,-8.45) {\textbf{Согласуйте единицы}\quad Запишите формульный столбец или относительный путь};
\node[flow] (e) at (0,-10.35) {\textbf{Составьте уравнение}\quad Свяжите времена, пути или скорости условием};
\node[flow] (f) at (0,-12.25) {\textbf{Решите и проверьте}\quad ОДЗ, смысл, единицы, именно искомая величина};
\draw[arr] (a)--(b);
\draw[arr] (b.west) -| node[pos=.28,above,font=\scriptsize]{река} (c1.north);
\draw[arr] (b.south)--node[right,font=\scriptsize]{средняя}(c2.north);
\draw[arr] (b.east)-|node[pos=.28,above,font=\scriptsize]{поезд}(c3.north);
\draw[draw=accent,line width=.85pt] (c1.south) |- (m);
\draw[draw=accent,line width=.85pt] (c2.south) -- (m);
\draw[draw=accent,line width=.85pt] (c3.south) |- (m);
\fill[accent] (m) circle (1.2pt);
\draw[arr] (m)--(d);
\draw[arr] (d)--(e);
\draw[arr] (e)--(f);
\end{tikzpicture}
\end{center}

\clearpage
\section*{07. Тренировочный блок · Путь А}
\ruleline
\textcolor{muted}{10 заданий по возрастанию сложности. Постройте математическую модель и запишите ответ с единицами, где они необходимы.}
\begin{enumerate}
\item Катер прошёл 24 км по течению и 24 км против течения. Скорость течения 3 км/ч. На путь против течения он затратил на 2 часа больше. Найдите собственную скорость катера.

\item Теплоход вышел в 7:00, прошёл 70 км по течению, простоял 2 часа и вернулся в исходный пункт в 21:00. Скорость течения 2 км/ч. Найдите собственную скорость теплохода.

\item Плот со скоростью 2 км/ч прошёл 22 км. Через час после отправления плота от того же пункта вышла моторная лодка, прошла 48 км по течению и 48 км против течения и вернулась одновременно с прибытием плота. Найдите собственную скорость лодки.

\item На одинаковые ненулевые расстояния баржа затрачивает по течению в 3 раза меньше времени, чем против течения. Во сколько раз собственная скорость баржи больше скорости течения?

\item Автомобиль половину общего времени двигался со скоростью 72 км/ч, а вторую половину --- со скоростью 48 км/ч. Найдите среднюю скорость за всю поездку.

\item Автомобиль половину всего пути двигался со скоростью 72 км/ч, а вторую половину --- со скоростью 48 км/ч. Найдите среднюю скорость за всю поездку.

\item Велосипедист прошёл первую треть всего пути со скоростью 30 км/ч, а оставшиеся две трети --- со скоростью 60 км/ч. Найдите среднюю скорость на всём пути.

\item Поезд движется со скоростью 90 км/ч и проходит мимо неподвижного столба за 24 секунды. Найдите длину поезда в метрах.

\item Поезд движется со скоростью 72 км/ч и полностью проходит мост длиной 900 м за 75 секунд. Найдите длину поезда в метрах.

\item Два поезда движутся по параллельным путям в одном направлении. Медленный поезд имеет длину 210 м и скорость 54 км/ч, быстрый --- длину 190 м. В начальный момент \textbf{голова} быстрого поезда находится в 200 м позади \textbf{хвоста} медленного. Через 60 секунд быстрый поезд целиком завершил обгон. Найдите скорость быстрого поезда в км/ч.
\end{enumerate}

\clearpage
\section*{08. Ответы для самопроверки}
\ruleline
\begin{center}
\renewcommand{\arraystretch}{1.65}
\begin{tabularx}{.90\linewidth}{>{\bfseries}cX}
\toprule
№ & Ответ\\
\midrule
1 & $9$ км/ч\\
2 & $12$ км/ч\\
3 & $10$ км/ч\\
4 & В $2$ раза\\
5 & $60$ км/ч\\
6 & $57{,}6$ км/ч\\
7 & $45$ км/ч\\
8 & $600$ м\\
9 & $600$ м\\
10 & $90$ км/ч\\
\bottomrule
\end{tabularx}
\end{center}
\vspace{6mm}
\key{Контрольные ориентиры.}
\begin{itemize}
\item В заданиях 1--3 собственная скорость строго больше скорости течения.
\item В заданиях 5 и 6 одинаковые скорости участков, но различны условия: равные времена и равные пути.
\item В заданиях 8--10 результат получен с согласованными единицами. В задании 10 относительный путь равен сумме длин поездов и начального зазора.
\end{itemize}
\end{document}
